\documentclass[a4paper,12pt]{article}
\usepackage[T2A]{fontenc}
\usepackage[utf8]{inputenc}
\usepackage[russian]{babel}
\usepackage{amsmath,amssymb}
\usepackage{PTSans}
\renewcommand{\familydefault}{\sfdefault}
\usepackage{icomma}
\usepackage[a4paper,margin=20mm,headheight=25pt,headsep=6mm]{geometry}
\usepackage{microtype}
\usepackage{tikz}
\usetikzlibrary{calc,arrows.meta,positioning,shapes.geometric}
\usepackage{tabularx,booktabs,longtable}
\usepackage{enumitem}
\usepackage{parskip}
\usepackage{fancyhdr}
\usepackage{setspace}
\usepackage{xcolor}

\definecolor{accent}{HTML}{008080}
\definecolor{darkaccent}{HTML}{004D4D}
\definecolor{textgray}{HTML}{333333}

\onehalfspacing
\setlength{\parindent}{0pt}
\setlist[itemize]{leftmargin=1.45em,itemsep=0.2em,topsep=0.25em}
\setlist[enumerate]{leftmargin=1.8em,itemsep=0.4em,topsep=0.3em}
\renewcommand{\arraystretch}{1.28}

\pagestyle{fancy}
\fancyhf{}
\fancyhead[L]{\small\textcolor{textgray}{Надя Клименко, 6 класс}}
\fancyhead[R]{\small\textcolor{textgray}{Математика}}
\fancyfoot[C]{\small\thepage}
\renewcommand{\headrulewidth}{0.4pt}
\renewcommand{\headrule}{%
  \hbox to\headwidth{%
    \color{accent}\leaders\hrule height \headrulewidth\hfill}}
\renewcommand{\footrulewidth}{0pt}

\newcommand{\sectionline}{%
  \par\vspace{0.1em}%
  {\color{accent}\hrule height 0.7pt}%
  \vspace{0.55em}}

\newcommand{\key}[1]{\textcolor{darkaccent}{\textbf{#1}}}
\newcommand{\GCDru}{\mathop{\text{НОД}}\nolimits}
\newcommand{\LCMru}{\mathop{\text{НОК}}\nolimits}

\tikzset{
  flowstart/.style={
    ellipse,
    draw=darkaccent,
    line width=0.8pt,
    align=center,
    minimum width=48mm,
    minimum height=12mm,
    font=\small,
    inner sep=4pt
  },
  flowaction/.style={
    rounded corners=3pt,
    draw=accent,
    line width=0.8pt,
    align=center,
    text width=74mm,
    minimum height=13mm,
    font=\small,
    inner sep=5pt
  },
  flowdecision/.style={
    diamond,
    aspect=2.15,
    draw=darkaccent,
    line width=0.8pt,
    align=center,
    text width=47mm,
    font=\small,
    inner sep=2pt
  },
  flowend/.style={
    ellipse,
    draw=darkaccent,
    line width=0.8pt,
    align=center,
    minimum width=52mm,
    minimum height=12mm,
    font=\small,
    inner sep=4pt
  },
  flowarrow/.style={
    -{Latex[length=3mm,width=2mm]},
    line width=0.8pt,
    draw=textgray
  }
}

\newenvironment{flowchartpage}{%
  \clearpage
  \thispagestyle{empty}%
}{%
  \clearpage
}

\newenvironment{flowchartpagelandscape}{%
  \clearpage
  \begingroup
  \pdfpagewidth=420mm
  \pdfpageheight=297mm
  \paperwidth=420mm
  \paperheight=297mm
  \newgeometry{left=22mm,right=22mm,top=20mm,bottom=20mm}
  \setlength{\textwidth}{376mm}%
  \setlength{\linewidth}{376mm}%
  \setlength{\hsize}{376mm}%
  \setlength{\columnwidth}{376mm}%
  \setlength{\textheight}{257mm}%
  \thispagestyle{empty}%
}{%
  \clearpage
  \restoregeometry
  \endgroup
  \pdfpagewidth=210mm
  \pdfpageheight=297mm
  \paperwidth=210mm
  \paperheight=297mm
}

\begin{document}

% -------------------- Титульный лист --------------------
\thispagestyle{empty}
\begin{center}

\vspace*{22mm}

{\Large\textcolor{darkaccent}{\textbf{Чек-лист}}}\par

\vspace{5mm}

{\LARGE\bfseries НОД и НОК: вычисление, смысл и связь}\par

\vspace{4mm}

{\Large\bfseries Простые уравнения и задачи на движение}\par

\vspace{18mm}

{\large\textbf{Надя Клименко, 6 класс}}\par

\vspace{8mm}

{\normalsize Тренировочный блок: \textbf{Путь А}}\par

\vfill

{\normalsize \today}\par

\vspace{5mm}

{\small Лёвин Артём Александрович эксклюзивно для Нади Клименко}\par

\vspace{12mm}

\begin{tikzpicture}[remember picture,overlay]
  \draw[accent!55,line width=1.1pt]
    ($(current page.south west)+(18mm,16mm)$)
      .. controls
      ($(current page.south west)+(70mm,28mm)$)
      and
      ($(current page.south east)+(-72mm,8mm)$)
      ..
      ($(current page.south east)+(-18mm,18mm)$);
\end{tikzpicture}

\end{center}

\clearpage

% -------------------- Теория --------------------
\section*{1. Простые уравнения: обратные действия}
\sectionline

Главная цель --- уверенно находить неизвестный компонент действия
и проверять ответ подстановкой.

\begin{tabularx}{\linewidth}{@{}>{\bfseries}l X X@{}}
\toprule
Тип & Как рассуждать & Пример \\
\midrule

$\dfrac{x}{a}=b$
&
Неизвестное делимое равно произведению делителя и частного.
&
$\dfrac{x}{3}=5 \Rightarrow x=15$.
\\

$\dfrac{x+c}{a}=b$
&
Сначала найдите $x+c$, затем $x$.
&
$\dfrac{x+1}{3}=5 \Rightarrow x=14$.
\\

$\dfrac{a}{x+c}=b$
&
Неизвестный делитель равен делимому, делённому на частное.
&
$\dfrac{3}{x+1}=5 \Rightarrow x=-\dfrac25$.
\\

$a+x=b$
&
Неизвестное слагаемое равно разности суммы и известного слагаемого.
&
$5+x=6 \Rightarrow x=1$.
\\

$a-x=b$
&
Вычитаемое равно разности уменьшаемого и разности.
&
$10-x=2 \Rightarrow x=8$.
\\

$ax+b=c$
&
Сначала уберите прибавление или вычитание,
затем делите на коэффициент при $x$.
&
$2x-3=10 \Rightarrow x=\dfrac{13}{2}$.
\\

\bottomrule
\end{tabularx}

\textbf{Ограничение.}
Если неизвестная находится в знаменателе,
знаменатель не может быть равен нулю.

Например,
\[
\frac{2}{x-3}=10,
\qquad
x\ne3.
\]

Так как
\[
x-3=\frac{2}{10}=\frac15,
\]
получаем
\[
x=\frac{16}{5}.
\]
Это значение допустимо.

\textbf{Самопроверка.}
Подставьте найденное число в исходное уравнение:
левая и правая части должны совпасть.

\section*{2. Разложение на простые множители}
\sectionline

Для вычисления НОД и НОК удобно раскладывать числа
на простые множители:
\[
84=2^2\cdot3\cdot7,
\qquad
126=2\cdot3^2\cdot7.
\]

Показатель степени показывает,
сколько раз соответствующий простой множитель
входит в разложение числа.

\textbf{Простое число} имеет ровно два натуральных делителя:
$1$ и само это число.
Число $1$ простым не является.

\section*{3. НОД: наибольший общий делитель}
\sectionline

\textbf{Определение.}
$\GCDru(a,b)$ --- наибольшее натуральное число,
на которое и $a$, и $b$ делятся без остатка.

Для НОД берём только
\key{общие простые множители}
в
\key{наименьших степенях}:
\[
\GCDru(84,126)
=
2^1\cdot3^1\cdot7
=
42.
\]

\textbf{Смысл в текстовых задачах.}
НОД особенно часто возникает,
когда нужно разбить несколько количеств
на одинаковые части без остатка
и сделать размер части или число групп
максимально возможным.

\textbf{Пример.}
Ленты длиной $84$ см и $126$ см
разрезают на одинаковые куски
наибольшей возможной длины.

Тогда
\[
\GCDru(84,126)=42,
\]
поэтому длина одного куска равна $42$ см.

Из первой ленты получится
\[
84:42=2
\]
куска, из второй ---
\[
126:42=3
\]
куска.

\section*{4. НОК: наименьшее общее кратное}
\sectionline

\textbf{Определение.}
$\LCMru(a,b)$ --- наименьшее натуральное число,
которое делится и на $a$, и на $b$.

Например,
\[
12=2^2\cdot3,
\qquad
18=2\cdot3^2.
\]

Для НОК берём
\key{все нужные простые множители}
в
\key{наибольших степенях}:
\[
\LCMru(12,18)
=
2^2\cdot3^2
=
36.
\]

\textbf{Смысл в текстовых задачах.}
НОК часто возникает при повторяющихся событиях:
требуется узнать,
через какое наименьшее время
несколько циклов снова совпадут.

\textbf{Пример.}
Один сигнал повторяется каждые $12$ минут,
другой --- каждые $18$ минут.
Если сейчас они прозвучали одновременно,
то снова одновременно прозвучат через
\[
\LCMru(12,18)=36\text{ минут}.
\]

% -------------------- Схема НОД / НОК --------------------
\begin{flowchartpagelandscape}

\begin{center}
{\Large\bfseries\textcolor{darkaccent}{Алгоритм: как выбрать НОД или НОК}}
\end{center}

\vspace{4mm}

\begin{center}
\begin{tikzpicture}[node distance=7mm]

  \node[flowstart] (start)
    {Определите, что именно требуется найти};

  \node[flowdecision,below=of start] (split)
    {Нужно разбить величины на одинаковые части без остатка?};

  \node[flowdecision,below=of split] (max)
    {Требуется максимально возможный размер части
     или максимум одинаковых групп?};

  \node[flowdecision,below=of max] (cycle)
    {Есть повторяющиеся циклы
     или требуется наименьшее общее кратное?};

  \node[flowaction,below=of cycle] (other)
    {Составьте другое уравнение по смыслу задачи};

  \node[flowend,below=of other] (finish)
    {Запишите ответ и объясните смысл найденного числа};

  \node[
    flowaction,
    left=13mm of cycle,
    text width=48mm
  ] (gcd)
    {Найдите НОД данных чисел};

  \node[
    flowaction,
    right=13mm of other,
    text width=48mm
  ] (lcm)
    {Найдите НОК данных чисел};

  \draw[flowarrow]
    (start) -- (split);

  \draw[flowarrow]
    (split)
    --
    node[right,font=\small]{да}
    (max);

  \draw[flowarrow]
    (split.east)
    --
    ++(8mm,0)
    |-
    node[pos=0.18,right,font=\small]{нет}
    (cycle.east);

  \draw[flowarrow]
    (max.west)
    --
    node[above,font=\small]{да}
    (gcd.north);

  \draw[flowarrow]
    (max)
    --
    node[right,font=\small]{нет}
    (cycle);

  \draw[flowarrow]
    (cycle.east)
    --
    node[above,font=\small]{да}
    (lcm.north);

  \draw[flowarrow]
    (cycle)
    --
    node[right,font=\small]{нет}
    (other);

  \draw[flowarrow]
    (gcd.south)
    |-
    (finish.west);

  \draw[flowarrow]
    (lcm.south)
    |-
    (finish.east);

  \draw[flowarrow]
    (other)
    --
    (finish);

\end{tikzpicture}
\end{center}

\end{flowchartpagelandscape}

% -------------------- Связь НОД и НОК --------------------
\section*{5. Связь НОД и НОК для двух чисел}
\sectionline

Для двух натуральных чисел $a$ и $b$ верна формула
\[
\boxed{
\GCDru(a,b)\cdot\LCMru(a,b)=ab
}.
\]

Отсюда
\[
\LCMru(a,b)
=
\frac{ab}{\GCDru(a,b)},
\qquad
\GCDru(a,b)
=
\frac{ab}{\LCMru(a,b)}.
\]

\textbf{Пример.}
Известно, что
\[
a=84,
\qquad
b=180,
\qquad
\GCDru(a,b)=12.
\]

Тогда
\[
\LCMru(a,b)
=
\frac{84\cdot180}{12}
=
84\cdot15
=
1260.
\]

\textbf{Важно.}
Формула
\[
\GCDru(a,b)\cdot\LCMru(a,b)=ab
\]
в таком виде относится к
\textbf{двум натуральным числам}.

Для трёх и более чисел нельзя
просто перемножить все числа
и разделить полученное произведение
на их общий НОД.

\subsection*{Обратный ход}

Если известны НОД и НОК двух чисел,
то известно и их произведение:
\[
ab
=
\GCDru(a,b)\cdot\LCMru(a,b).
\]

При поиске самих чисел
возможные пары нужно дополнительно проверять
по НОД и НОК.

\section*{6. Степени простых множителей}
\sectionline

Если
\[
n=2^x3^y5^z7^w,
\]
то $x,y,z,w$ показывают,
сколько раз соответствующие простые множители
входят в разложение числа $n$.

\textbf{Пример.}
Пусть
\[
\GCDru(n,360)=120,
\qquad
\LCMru(n,360)=2520.
\]

Для двух чисел $n$ и $360$:
\[
120\cdot2520
=
n\cdot360.
\]

Следовательно,
\[
n
=
\frac{120\cdot2520}{360}
=
\frac{2520}{3}
=
840.
\]

Разложим число $840$:
\[
840
=
2^3\cdot3\cdot5\cdot7.
\]

Значит,
\[
x=3,
\qquad
y=z=w=1.
\]

Полезно помнить общий принцип.
Если
\[
a=2^{\alpha}3^{\beta},
\qquad
b=2^{\gamma}3^{\delta},
\]
то
\[
\GCDru(a,b)
=
2^{\min(\alpha,\gamma)}
3^{\min(\beta,\delta)},
\]
а
\[
\LCMru(a,b)
=
2^{\max(\alpha,\gamma)}
3^{\max(\beta,\delta)}.
\]

\section*{7. Простое число в задаче на НОК}
\sectionline

\textbf{Пример.}
Найдите простое число $p$, если
\[
\LCMru(p,18)=90.
\]

Имеем
\[
18=2\cdot3^2,
\qquad
90=2\cdot3^2\cdot5.
\]

Чтобы НОК увеличился от $18$ до $90$,
должен появиться новый простой множитель $5$.
Поэтому
\[
p=5.
\]

Числа $2$ и $3$ уже входят
в разложение числа $18$,
поэтому при $p=2$ или $p=3$
НОК остался бы равен $18$.

\section*{8. Задачи на движение: от текста к уравнению}
\sectionline

Для равномерного движения используется формула
\[
S=vt,
\]
где $S$ --- путь,
$v$ --- скорость,
$t$ --- время.

\textbf{Пример.}
Из двух пристаней навстречу друг другу
одновременно вышли катер и теплоход.
Расстояние между пристанями $58$ км.
Скорость теплохода на $2$ км/ч
больше скорости катера.
Они встретились через $2$ часа.
Найдите обе скорости.

Пусть скорость катера равна $x$ км/ч.
Тогда скорость теплохода равна
$x+2$ км/ч.

\begin{table}[ht]
\centering
\caption{Таблица величин для задачи на движение}
\begin{tabularx}{0.92\linewidth}{@{}lXXX@{}}
\toprule
Участник
&
Скорость, км/ч
&
Время, ч
&
Путь, км
\\
\midrule
Катер
&
$x$
&
$2$
&
$2x$
\\
Теплоход
&
$x+2$
&
$2$
&
$2(x+2)$
\\
\bottomrule
\end{tabularx}
\end{table}

До встречи катер и теплоход
вместе прошли всё расстояние между пристанями:
\[
2x+2(x+2)=58.
\]

Решаем:
\[
2x+2x+4=58,
\]
\[
4x=54,
\]
\[
x=13{,}5.
\]

Следовательно, скорости равны
\[
13{,}5\text{ км/ч}
\qquad\text{и}\qquad
15{,}5\text{ км/ч}.
\]

Проверка:
\[
2\cdot13{,}5
+
2\cdot15{,}5
=
27+31
=
58.
\]

% -------------------- Схема движения --------------------
\begin{flowchartpage}

\begin{center}
{\Large\bfseries\textcolor{darkaccent}{Алгоритм: задача на движение навстречу}}
\end{center}

\vspace{4mm}

\begin{center}
\begin{tikzpicture}[node distance=6.5mm]

  \node[flowstart] (start)
    {Выпишите известные величины и то, что нужно найти};

  \node[flowaction,below=of start] (var)
    {Обозначьте неизвестную скорость через $x$
     и выразите остальные скорости};

  \node[flowaction,below=of var] (table)
    {Заполните таблицу $v$, $t$, $S$;
     путь вычисляйте по формуле $S=vt$};

  \node[flowdecision,below=of table] (meet)
    {Участники движутся навстречу
     и встречаются между пунктами?};

  \node[flowaction,below=of meet] (other)
    {Составьте уравнение по конкретному условию движения};

  \node[flowaction,below=of other] (sum)
    {Для встречи навстречу:
     сложите пути и приравняйте сумму
     к общему расстоянию};

  \node[flowaction,below=of sum] (solve)
    {Решите полученное уравнение
     и найдите требуемые величины};

  \node[flowdecision,below=of solve] (check)
    {Подстановка и смысл ответа
     подтверждают результат?};

  \node[flowaction,below=of check] (repair)
    {Исправьте таблицу, уравнение или вычисления};

  \node[flowend,below=of repair] (finish)
    {Запишите ответ с единицами измерения};

  \draw[flowarrow]
    (start)
    --
    (var);

  \draw[flowarrow]
    (var)
    --
    (table);

  \draw[flowarrow]
    (table)
    --
    (meet);

  \draw[flowarrow]
    (meet)
    --
    node[right,font=\small]{нет}
    (other);

  \draw[flowarrow]
    (meet.west)
    --
    ++(-9mm,0)
    |-
    node[pos=0.18,left,font=\small]{да}
    (sum.west);

  \draw[flowarrow]
    (other.east)
    --
    ++(9mm,0)
    |-
    (solve.east);

  \draw[flowarrow]
    (sum)
    --
    (solve);

  \draw[flowarrow]
    (solve)
    --
    (check);

  \draw[flowarrow]
    (check)
    --
    node[right,font=\small]{нет}
    (repair);

  \draw[flowarrow]
    (check.west)
    --
    ++(-9mm,0)
    |-
    node[pos=0.18,left,font=\small]{да}
    (finish.west);

  \draw[flowarrow]
    (repair.east)
    --
    ++(11mm,0)
    |-
    (table.east);

\end{tikzpicture}
\end{center}

\end{flowchartpage}

% -------------------- Ошибки и итог --------------------
\section*{9. Типичные ошибки}
\sectionline

\begin{enumerate}

  \item
  \textbf{Путать НОД и НОК.}
  Разбиение на одинаковые максимально крупные части
  обычно связано с НОД;
  совпадение циклов и наименьшее общее кратное ---
  с НОК.

  \item
  \textbf{Неверно выбирать степени.}
  В НОД берутся меньшие показатели
  общих простых множителей;
  в НОК --- большие показатели
  всех нужных простых множителей.

  \item
  \textbf{Применять формулу двух чисел
  к трём и более числам.}
  Равенство
  \[
  \GCDru(a,b)\LCMru(a,b)=ab
  \]
  напрямую работает только для пары чисел.

  \item
  \textbf{Забывать ограничение в дробном уравнении.}
  Знаменатель не может быть равен нулю.

  \item
  \textbf{Получать число без объяснения его смысла.}
  В текстовой задаче нужно указать,
  что означает ответ:
  длину куска, число минут, скорость и т.~д.

  \item
  \textbf{Составлять уравнение на движение
  без таблицы и проверки.}
  Для первых задач надёжнее сначала выразить величины,
  заполнить таблицу,
  затем составить уравнение.

\end{enumerate}

\section*{10. Чек-лист перед самостоятельной работой}
\sectionline

Убедитесь, что Вы умеете:

\begin{itemize}

  \item
  решать простые линейные и дробные уравнения
  и проверять ответ;

  \item
  раскладывать натуральные числа
  на простые множители;

  \item
  находить НОД и НОК по разложению;

  \item
  выбирать НОД или НОК
  по смыслу текстовой задачи;

  \item
  применять формулу
  \[
  \GCDru(a,b)\cdot\LCMru(a,b)=ab
  \]
  только для двух чисел;

  \item
  восстанавливать число
  по известным НОД и НОК;

  \item
  читать показатели степеней
  в разложении на простые множители;

  \item
  использовать простоту числа
  в задачах на НОК;

  \item
  составлять таблицу движения
  и переводить условие в уравнение.

\end{itemize}

% -------------------- Тренировка --------------------
\clearpage

\section*{Тренировочный блок --- Путь А}
\sectionline

\begin{enumerate}[label=\textbf{\arabic*.}]

  \item
  Решите уравнение
  \[
  \frac{8}{x-2}=4.
  \]
  Сначала укажите значение $x$,
  при котором выражение не имеет смысла.
  После решения выполните проверку подстановкой.

  \item
  Известно, что для двух натуральных чисел
  $a=168$ и $b$ выполняется
  \[
  \GCDru(a,b)=24,
  \qquad
  \LCMru(a,b)=840.
  \]
  Найдите число $b$
  и разложите его на простые множители.

  \item
  Из двух пунктов,
  расстояние между которыми $82$ км,
  одновременно навстречу друг другу
  выехали два велосипедиста.
  Скорость второго на $5$ км/ч
  больше скорости первого.
  Через $2$ часа они встретились.
  Найдите скорости обоих велосипедистов.

\end{enumerate}

% -------------------- Ответы --------------------
\clearpage

\section*{Ответы к тренировочному блоку}
\sectionline

\begin{longtable}{@{}p{0.12\linewidth}p{0.82\linewidth}@{}}

\toprule
\textbf{№} & \textbf{Ответ} \\
\midrule
\endfirsthead

\toprule
\textbf{№} & \textbf{Ответ} \\
\midrule
\endhead

1
&
Ограничение:
$x\ne2$.
Ответ:
$x=4$.
Проверка:
\[
\frac{8}{4-2}=4.
\]
\\

2
&
\[
b=120,
\qquad
120=2^3\cdot3\cdot5.
\]
\\

3
&
$18$ км/ч и $23$ км/ч.
\\

\bottomrule

\end{longtable}

\clearpage

\end{document}